\section{Reproduction des campagnes}
\label{annexe:reproduction}

Toutes les commandes s'ex\'ecutent depuis la racine du d\'ep\^ot. La variable
\code{\$m} d\'esigne le mod\`ele de r\'ef\'erence :

\begin{lstlisting}[language=bash, caption={Variables et construction.}]
$m = 'python_src/checkpoints_onnx/2026_04_23_23h25_iter316_unsupervised.onnx'
uv run pytest python_src/tests -q
ctest --test-dir build -C Release --output-on-failure
\end{lstlisting}

\subsection*{D\'ebit et phases}

\begin{lstlisting}[language=bash, caption={\'Etape 10 du plan, campagne
principale et campagne de phases.}]
uv run python python_src/search_bench.py --model $m --gpu \
  --simulations 700 --passages 5 --repetitions 6 --batch-sizes 8 \
  --worker-counts 1 2 4 8 --cache-history-depths 0 --tt-size 8192 \
  --warmup-pool --out-json out/multicore/after.json \
  --out-rapport out/multicore/after.md

uv run python python_src/search_bench.py --model $m --gpu \
  --simulations 700 --passages 1 --repetitions 3 --batch-sizes 8 \
  --worker-counts 1 2 4 8 --cache-history-depths 0 --timings \
  --warmup-pool --out-json out/multicore/phases.json \
  --out-rapport out/multicore/phases.md
\end{lstlisting}

La campagne de queue utilise les m\^emes options avec
\code{--passages 50 --repetitions 1 --worker-counts 1 2 8} et
\code{--out-json out/multicore/queue.json}.

\subsection*{Qualit\'e}

\begin{lstlisting}[language=bash, caption={Pr\'efiltre, campagne compl\`ete et
temps \'egal.}]
# prefiltre 500 puzzles
uv run python python_src/puzzle_bench.py --model $m \
  --banc data/puzzles_bench.txt --gpu --travailleurs 1 \
  --search-workers 1 --batch-size 8 --cache-history-depth 0 \
  --simulations 700 --limite 500 \
  --out-csv out/multicore/prefilter-w1.csv

# campagne complete 2500 (idem pour --search-workers 8)
uv run python python_src/puzzle_bench.py --model $m \
  --banc data/puzzles_bench.txt --gpu --travailleurs 1 \
  --search-workers 8 --batch-size 8 --cache-history-depth 0 \
  --simulations 700 --limite 2500 \
  --out-csv out/multicore/full-w8.csv

# temps egal, budget median observe de 700 simulations mono-worker
uv run python python_src/puzzle_bench.py --model $m \
  --banc data/puzzles_bench.txt --gpu --travailleurs 1 \
  --search-workers 8 --batch-size 8 --cache-history-depth 0 \
  --search-seconds 1.5 --limite 500 \
  --out-csv out/multicore/temps-w8.csv

# comparaison appariee (sidecar obligatoire a cote de chaque CSV)
uv run python python_src/multicore_comparison.py \
  --baseline out/multicore/full-w1.csv \
  --candidate out/multicore/full-w8.csv \
  --out-report out/multicore/full-comparison.md --seed 42
\end{lstlisting}

\subsection*{Self-play partag\'e}

\begin{lstlisting}[language=bash, caption={Banc self-play. Le dossier
\code{torch/lib} doit \^etre dans \code{PATH} pour que le provider CUDA trouve
\code{cublasLt}.}]
build/Release/selfplay_phase_bench.exe --fake \
  > out/multicore/selfplay-after-fake.json

$env:PATH = ".venv\Lib\site-packages\torch\lib;$env:PATH"
build/Release/selfplay_phase_bench.exe --model $m --gpu \
  > out/multicore/selfplay-after-gpu.json
\end{lstlisting}

\subsection*{Positions de r\'ef\'erence}

\begin{center}
\small
\begin{tabular}{llp{8.2cm}}
\toprule
Nom & R\^ole & FEN \\
\midrule
ouverture & attaque italienne & \code{r1bqkbnr/1ppp1ppp/p1n5/1B2p3/4P3/} \\
 & & \code{5N2/PPPP1PPP/RNBQK2R w KQkq - 0 4} \\
milieu & position de Kiwipete & \code{r3k2r/p1ppqpb1/bn2pnp1/3PN3/1p2P3/} \\
 & & \code{2N2Q1p/PPPBBPPP/R3K2R w KQkq - 0 1} \\
finale & finale de tours & \code{8/2p5/3p4/KP5r/1R3p1k/8/4P1P1/8 w - - 0 1} \\
\bottomrule
\end{tabular}
\end{center}
